%
% FH Technikum Wien
% !TEX encoding = UTF-8 Unicode
%
% Version V2.24 von 2024-12-19 otrebski
%
\documentclass[BMR,Seminar,ngerman,IEEE]{twbook}
\renewcommand*{\degreecourse}{FH Technikum Wien, Bachelorstudiengang BIF\\Lehrveranstaltung BIF-VZ-3-WS2025-CLCOM-EN/158471}

% --- Encoding & fonts ---
\usepackage[utf8]{inputenc}
\usepackage[T1]{fontenc}

% --- Graphics, figures, urls, lists, tables, code ---
\usepackage{float}
\usepackage{placeins}
\usepackage{graphicx}
\usepackage{caption}
\usepackage{booktabs}
\usepackage{tabularx}
\usepackage{enumitem}
\usepackage{listings}
\usepackage[final]{microtype} % hilft gegen Overfull/Underfull
\emergencystretch=2em         % sanfter Puffer für Zeilenumbruch

% Screenshots hier hinein legen:
\graphicspath{{screenshots/}}

% --- Bibliography (twbook lädt biblatex via Hook, sobald inputenc geladen ist) ---
\addbibresource{Literatur.bib}

% --- Listings (inkl. YAML) ---
\lstdefinelanguage{yaml}{
  keywords={true,false,null,y,n},
  comment=[l]{\#},
  morestring=[b]',
  morestring=[b]",
  sensitive=false,
  showstringspaces=false,
  literate =    {---}{{\textcolor{gray}{---}}}3
                {>}{{\textcolor{gray}{>}}}1
                {|}{{\textcolor{gray}{|}}}1
                {:}{{\textcolor{blue}{:}}}1
                {-}{{\textcolor{gray}{-}}}1,
}
\lstset{
  basicstyle=\ttfamily\small,
  frame=single,
  breaklines=true,
  backgroundcolor=\color{gray!5},
  captionpos=b
}

% --- Convenience figure macros ---
\newcommand{\screenshotH}[3]{%
  \begin{figure}[H]
    \centering
    \includegraphics[width=\linewidth]{#1}%
    \caption{#2}%
    \label{fig:#3}%
  \end{figure}%
}
\newcommand{\placeholderfig}[2]{%
  \begin{figure}[htbp]
    \centering
    \fbox{\rule{0pt}{120pt}\rule{0.9\linewidth}{0pt}}%
    \caption{#1}%
    \label{fig:#2}%
  \end{figure}%
}

% Die nachfolgenden Pakete stellen sonst nicht benötigte Features zur Verfügung
\usepackage{blindtext}

% --- Hyperref am Schluss laden (robuster) ---
\usepackage{hyperref}

%
% Einträge für Deckblatt, Kurzfassung, etc.
%
\title{Semesterprojekt – DevOps und Cloud Computing 3B}
\author{Benjamin Feichtlbauer, Kevin Forter, Mikulovic Luka \and Tepsurkaev Abdulah}
\place{Wien}
% \acknowledgements{\blindtext}

\begin{document}

\maketitle

% ============================================================================
\onecolumn
\newpage

% ============================================================================
\chapter{Milestone 2}

\section{Network Setup}
In this section, we describe the setup of the Virtual Private Cloud (VPC) and the associated network components within AWS.

\subsection{VPC Configuration}
We created a new VPC with the IPv4 CIDR block \texttt{10.0.0.0/16}. This provides a private network space for our infrastructure.

\screenshotH{assets/setup/vpc.png}{VPC Configuration}{vpc}

\subsection{Subnets}
A public subnet was created within the VPC with the CIDR block \texttt{10.0.0.0/24}. This subnet hosts our instances and allows them to communicate with each other and the internet.

\screenshotH{assets/setup/subnet.png}{Subnet Configuration}{subnet}
\screenshotH{assets/setup/subnet_settings.png}{Subnet Settings}{subnet_settings}

\subsection{Internet Gateway}
To enable internet access for the resources in our public subnet, we attached an Internet Gateway to the VPC.

\screenshotH{assets/setup/igw.png}{Internet Gateway}{igw}

\subsection{Route Tables}
A custom route table was configured to direct traffic destined for the internet (\texttt{0.0.0.0/0}) to the Internet Gateway. This route table is associated with our public subnet.

\screenshotH{assets/setup/routetable.png}{Route Table}{routetable}
\screenshotH{assets/setup/routetable_default_route.png}{Route Table Default Route}{routetable_route}
\screenshotH{assets/setup/routetable_subnet_associations.png}{Route Table Subnet Association}{routetable_association}

\subsection{Security Groups}
Specific security groups were created for each component to control inbound and outbound traffic.

\screenshotH{assets/setup/dns_security_group.png}{DNS Security Group}{sg_dns}
\screenshotH{assets/setup/gitlab_security_group.png}{GitLab Security Group}{sg_gitlab}
\screenshotH{assets/setup/runnger_security_group.png}{Runner Security Group}{sg_runner}

\section{DNS Setup}
We implemented a split-horizon DNS architecture using Bind9, with a primary and a secondary DNS server to ensure redundancy. The domain configured is \texttt{ci.intern}.

\subsection{Primary DNS}
The primary DNS server is hosted on an EC2 instance with the private IP \texttt{10.0.0.10}.

\screenshotH{assets/instances/dns1/details.png}{Primary DNS Instance Details}{dns1_details}

\subsubsection{Installation}
We installed Bind9 on the Ubuntu 22.04 LTS instance:
\begin{lstlisting}[language=bash]
sudo apt update
sudo apt install bind9 bind9utils -y
\end{lstlisting}

\subsubsection{Configuration}
The zone configuration in \texttt{/etc/bind/named.conf.local} defines the forward and reverse lookup zones:

\begin{lstlisting}
zone "ci.intern" {
    type master;
    file "/etc/bind/db.ci.intern";
    allow-transfer { 10.0.0.11; };
};

zone "0.0.10.in-addr.arpa" {
    type master;
    file "/etc/bind/db.10";
    allow-transfer { 10.0.0.11; };
};
\end{lstlisting}

The forward zone file \texttt{/etc/bind/db.ci.intern} maps hostnames to IP addresses:

\begin{lstlisting}
$TTL 3600
@   IN  SOA primary.dns.ci.intern. admin.ci.intern. (
        2025010101 ; Serial
        3600       ; Refresh
        1800       ; Retry
        604800     ; Expire
        86400 )    ; Negative Cache TTL

    IN  NS  primary.dns.ci.intern.
    IN  NS  secondary.dns.ci.intern.

primary.dns     IN A 10.0.0.10
secondary.dns   IN A 10.0.0.11
gitlab          IN A 10.0.0.20
runner          IN A 10.0.0.21
\end{lstlisting}

The reverse zone file \texttt{/etc/bind/db.10} maps IP addresses back to hostnames:

\begin{lstlisting}
$TTL 3600
@   IN  SOA primary.dns.ci.intern. admin.ci.intern. (
        2025010101
        3600
        1800
        604800
        86400 )

    IN NS primary.dns.ci.intern.
    IN NS secondary.dns.ci.intern.

10  IN PTR primary.dns.ci.intern.
11  IN PTR secondary.dns.ci.intern.
20  IN PTR gitlab.ci.intern.
21  IN PTR runner.ci.intern.
\end{lstlisting}

\subsection{Secondary DNS}
The secondary DNS server is hosted on an EC2 instance with the private IP \texttt{10.0.0.11}. It acts as a slave to the primary server.

\screenshotH{assets/instances/dns2/details.png}{Secondary DNS Instance Details}{dns2_details}

\subsubsection{Configuration}
The configuration in \texttt{/etc/bind/named.conf.local} specifies the master server for zone transfers:

\begin{lstlisting}
zone "ci.intern" {
    type slave;
    masters { 10.0.0.10; };
    file "/var/cache/bind/db.ci.intern";
};

zone "0.0.10.in-addr.arpa" {
    type slave;
    masters { 10.0.0.10; };
    file "/var/cache/bind/db.10";
};
\end{lstlisting}

\section{GitLab Server Setup}
We deployed a self-hosted GitLab CE server using Docker on an EC2 instance with the private IP \texttt{10.0.0.20}.

\screenshotH{assets/instances/gitlab/details.png}{GitLab Server Instance Details}{gitlab_details}

\subsection{Installation}
Docker and Docker Compose were installed, and a persistent volume structure was created. The \texttt{docker-compose.yml} configuration is as follows:

\begin{lstlisting}[language=yaml]
version: '3.7'
services:
  gitlab:
    image: gitlab/gitlab-ce:latest
    hostname: "gitlab.ci.intern"
    restart: always
    container_name: gitlab
    ports:
      - "80:80"
      - "443:443"
    volumes:
      - /srv/gitlab/config:/etc/gitlab
      - /srv/gitlab/logs:/var/log/gitlab
      - /srv/gitlab/data:/var/opt/gitlab
    environment:
      GITLAB_OMNIBUS_CONFIG: |
        external_url 'http://gitlab.ci.intern'
\end{lstlisting}

\subsection{Verification}
The GitLab instance is accessible via the browser using the public IP or the internal DNS name (within the network).

\screenshotH{assets/instances/gitlab/homepage.png}{GitLab Homepage}{gitlab_homepage}
\screenshotH{assets/instances/gitlab/login.png}{GitLab Login}{gitlab_login}

\section{GitLab Runner Setup}
A GitLab Runner was set up on a separate EC2 instance (\texttt{10.0.0.21}) to execute CI/CD jobs.

\screenshotH{assets/instances/runner/details.png}{GitLab Runner Instance Details}{runner_details}

\subsection{Configuration}
The runner was configured to use the internal DNS servers to resolve the GitLab server's hostname.

\begin{lstlisting}[language=bash]
sudo systemctl disable --now systemd-resolved
sudo rm /etc/resolv.conf
echo "nameserver 10.0.0.10" | sudo tee /etc/resolv.conf
echo "nameserver 10.0.0.11" | sudo tee -a /etc/resolv.conf
\end{lstlisting}

The runner container was started with specific DNS settings:

\begin{lstlisting}[language=bash]
sudo docker run -d --name gitlab-runner \
  --restart always \
  -v /srv/gitlab-runner/config:/etc/gitlab-runner \
  --dns 10.0.0.10 \
  --dns 10.0.0.11 \
  gitlab/gitlab-runner:latest
\end{lstlisting}

\subsection{Registration}
The runner was registered with the GitLab server using the internal URL \texttt{http://gitlab.ci.intern}.

\begin{lstlisting}[language=bash]
sudo docker exec -it gitlab-runner gitlab-runner register
# URL: http://gitlab.ci.intern
# Executor: docker
# Default Docker image: alpine:latest
\end{lstlisting}

\end{document}
